\chapter{2006年辽宁卷（文）}
\section{单选题}
\begin{problem}
函数 \( y = \sin \left(\frac{1}{2} x + 3\right) \) 的最小正周期是（\quad）

\choices
    {\( \frac{\pi}{2} \)}
    {\(\pi\)}
    {\(2\pi\)}
    {\(4\pi\)}
\end{problem}

\begin{answer}
D
\end{answer}

\begin{solution}
正弦函数 \(y=\sin(\omega x+\varphi)\) 的最小正周期为 \(T=\dfrac{2\pi}{|\omega|}\). 本题中 \(\omega=\dfrac{1}{2}\)，所以 \(T=\dfrac{2\pi}{1/2}=4\pi\)，故选 D.
\end{solution}

\begin{problem}
设集合 \(A = \{1, 2\}\)，则满足 \(A \cup B = \{1, 2, 3\}\) 的集合 \(B\) 的个数是（\quad）

\choices
    {\(1\)}
    {\(3\)}
    {\(4\)}
    {\(8\)}
\end{problem}

\begin{answer}
C
\end{answer}

\begin{solution}
由 \(A\cup B=\{1,2,3\}\) 可知 \(3\in B\)，且 \(B\) 中不能含有 \(1\)，\(2\)，\(3\) 以外的元素；元素 \(1,2\) 是否属于 \(B\) 均不影响并集. 因此，\(1,2\) 各有“取或不取”两种选择，集合 \(B\) 的个数为 \(2^2=4\)，故选 C.
\end{solution}

\begin{problem}
设 \( f(x) \) 是 \( \R \) 上的任意函数，则下列叙述正确的是（\quad）

\choices
    {\(f(x)f(-x)\) 是奇函数}
    {\(f(x)|f(-x)|\) 是奇函数}
    {\(f(x) + f(-x)\) 是偶函数}
    {\(f(x) - f(-x)\) 是偶函数}
\end{problem}

\begin{answer}
C
\end{answer}

\begin{solution}
令 \(F(x)=f(x)+f(-x)\)，则 \(F(-x)=f(-x)+f(x)=F(x)\)，所以 \(F(x)\) 是偶函数.

其余叙述不恒成立：取 \(f(x)=1\)，则 \(f(x)f(-x)=1\) 和 \(f(x)|f(-x)|=1\) 都不是奇函数；取 \(f(x)=x\)，则 \(f(x)-f(-x)=2x\) 是奇函数而不是偶函数. 因此只有第三个叙述正确，故选 C.
\end{solution}

\begin{problem}
\(\mathrm{C}_6^1 + \mathrm{C}_6^2 + \mathrm{C}_6^3 + \mathrm{C}_6^4 + \mathrm{C}_6^5\) 的值为（\quad）

\choices
    {\(61\)}
    {\(62\)}
    {\(63\)}
    {\(64\)}
\end{problem}

\begin{answer}
B
\end{answer}

\begin{solution}
由二项式系数和 \(\sum_{k=0}^{6}\mathrm C_6^k=2^6=64\)，得 \(\mathrm C_6^1+\mathrm C_6^2+\mathrm C_6^3+\mathrm C_6^4+\mathrm C_6^5=64-\mathrm C_6^0-\mathrm C_6^6=64-1-1=62\)，故选 B.
\end{solution}

\begin{problem}
方程 \(2x^{2} - 5x + 2 = 0\) 的两个根可分别作为（\quad）

\choices
    {一个椭圆和一双曲线的离心率}
    {两抛物线的离心率}
    {一个椭圆和一抛物线的离心率}
    {两椭圆的离心率}
\end{problem}

\begin{answer}
A
\end{answer}

\begin{solution}
将方程因式分解得 \((2x-1)(x-2)=0\)，所以两个根为 \(x_1=\dfrac{1}{2}\)，\(x_2=2\). 椭圆的离心率满足 \(0<e<1\)，双曲线的离心率满足 \(e>1\)，抛物线的离心率等于 \(1\). 因此 \(\dfrac{1}{2}\) 和 \(2\) 可分别作为一个椭圆和一条双曲线的离心率，故选 A.
\end{solution}

\begin{problem}
给出下列四个命题：

\begin{circlelist}
\item 垂直于同一直线的两条直线互相平行；
\item 垂直于同一平面的两个平面互相平行；
\item 若直线 \(l_{1}\)，\(l_{2}\) 与同一平面所成的角相等，则 \(l_{1}\)，\(l_{2}\) 互相平行；
\item 若直线 \(l_{1}\)，\(l_{2}\) 是异面直线，则与 \(l_{1}\)，\(l_{2}\) 都相交的两条直线是异面直线.
\end{circlelist}

其中假命题的个数是（\quad）

\choices
    {\(1\)}
    {\(2\)}
    {\(3\)}
    {\(4\)}
\end{problem}

\begin{answer}
D
\end{answer}

\begin{solution}
命题\circled{1}不正确：取同一直线 \(l\) 上一点 \(O\)，过 \(O\) 作两条不同的直线 \(l_1\)、\(l_2\)，使 \(l_1\perp l\)、\(l_2\perp l\)，则 \(l_1\)、\(l_2\) 相交而不平行.

命题\circled{2}不正确：取一个平面 \(\alpha\)，作两个都垂直于 \(\alpha\) 且互相垂直的平面 \(\beta\)、\(\gamma\)，则 \(\beta\)、\(\gamma\) 都垂直于 \(\alpha\)，但彼此不平行.

命题\circled{3}不正确：取同一平面内两条相交直线，它们与该平面所成的角都为 \(0\)，但两条直线不平行.

命题\circled{4}不正确：在 \(l_1\) 上取点 \(P\)，在 \(l_2\) 上取两个不同的点 \(Q_1\)、\(Q_2\)，则直线 \(PQ_1\)、\(PQ_2\) 都分别与 \(l_1\)、\(l_2\) 相交，但它们在 \(P\) 处相交而不是异面直线. 因此四个命题都是假命题，假命题的个数为 \(4\)，故选 D.
\end{solution}

\begin{problem}
双曲线 \(x^{2} - y^{2} = 4\) 的两条渐近线与直线 \(x = 3\) 围成一个三角形区域，表示该区域的不等式组是（\quad）

\choices
    {\(\begin{cases}
x - y \geqslant 0,\\
x + y \geqslant 0,\\
0 \leqslant x \leqslant 3.
\end{cases}\)}
    {\(\begin{cases}
x - y \geqslant 0,\\
x + y \leqslant 0,\\
0 \leqslant x \leqslant 3.
\end{cases}\)}
    {\(\begin{cases}
x - y \leqslant 0,\\
x + y \leqslant 0,\\
0 \leqslant x \leqslant 3.
\end{cases}\)}
    {\(\begin{cases}
x - y \leqslant 0,\\
x + y \geqslant 3,\\
0 \leqslant x \leqslant 3.
\end{cases}\)}
\end{problem}

\begin{answer}
A
\end{answer}

\begin{solution}
由 \(x^2-y^2=4\) 可知双曲线的两条渐近线为 \(y=x\) 和 \(y=-x\)，即 \(x-y=0\) 与 \(x+y=0\). 直线 \(x=3\) 与两条渐近线围成的三角形位于右侧夹角内，所以其中的点满足 \(y\leqslant x\)、\(y\geqslant-x\)，即 \(x-y\geqslant0\)、\(x+y\geqslant0\)，并且 \(0\leqslant x\leqslant3\). 故选 A.
\end{solution}

\begin{problem}
设 \(\oplus\) 是 \(\R\) 上的一个运算，\(A\) 是 \(\R\) 的非空子集，若对任意 \(a\)，\(b \in A\)，有 \(a \oplus b \in A\)，则称 \(A\) 对运算 \(\oplus\) 封闭. 下列数集对加法，减法，乘法和除法 （除数不等于零） 四则运算都封闭的是（\quad）

\choices
    {自然数集}
    {整数集}
    {有理数集}
    {无理数集}
\end{problem}

\begin{answer}
C
\end{answer}

\begin{solution}
自然数集对减法和除法不封闭，例如 \(1-2=-1\notin\N\)、\(1\div2=\dfrac12\notin\N\)；整数集对除法不封闭，例如 \(1\div2=\dfrac12\notin\Z\). 无理数集对加法也不封闭，例如 \(\sqrt2+(-\sqrt2)=0\) 是有理数. 任意两个有理数的和、差、积仍为有理数，除数不为零时商也仍为有理数，所以有理数集对题述四则运算都封闭，故选 C.
\end{solution}

\begin{problem}
\(\triangle ABC\) 的三内角 \(A\)，\(B\)，\(C\)，所对边的长分别为 \(a\)，\(b\)，\(c\)，设向量 \(\bs{p} = (a + c, b)\)，\(\bs{q} = (b - a, c - a)\). 若 \(\bs{p} \parallel \bs{q}\)，则角 \(C\) 的大小为（\quad）

\choices
    {\(\frac{\pi}{6}\)}
    {\(\frac{\pi}{3}\)}
    {\(\frac{\pi}{2}\)}
    {\(\frac{2\pi}{3}\)}
\end{problem}

\begin{answer}
B
\end{answer}

\begin{solution}
向量平行的充要条件给出 \((a+c)(c-a)-b(b-a)=0\). 化简得 \(c^2-a^2-b^2+ab=0\)，即 \(c^2=a^2+b^2-ab\). 另一方面，由余弦定理，\(c^2=a^2+b^2-2ab\cos C\)，所以 \(2ab\cos C=ab\). 由于三角形边长 \(a\)，\(b>0\)，得 \(\cos C=\dfrac12\)，且 \(0<C<\pi\)，从而 \(C=\dfrac{\pi}{3}\)，故选 B.
\end{solution}

\begin{problem}
已知等腰 \(\triangle ABC\) 的腰为底的 2 倍，则顶角 \(A\) 的正切值是（\quad）

\choices
    {\(\frac{\sqrt{3}}{2}\)}
    {\(\sqrt{3}\)}
    {\(\frac{\sqrt{15}}{8}\)}
    {\(\frac{\sqrt{15}}{7}\)}
\end{problem}

\begin{answer}
D
\end{answer}

\begin{solution}
设等腰三角形的底边长为 \(b\)，两腰长为 \(2b\). 由余弦定理，在顶角 \(A\) 对面有 \(b^2=(2b)^2+(2b)^2-2(2b)(2b)\cos A\)，所以 \(\cos A=\dfrac78\). 由于 \(0<A<\pi\)，\(\sin A=\sqrt{1-\dfrac{49}{64}}=\dfrac{\sqrt{15}}8\)，故 \(\tan A=\dfrac{\sin A}{\cos A}=\dfrac{\sqrt{15}}7\)，故选 D.
\end{solution}

\begin{problem}
与方程 \( y = \e^{2x} - 2\e^{x} + 1 \)（\(x \geqslant 0\)）的曲线关于直线 \( y = x \) 对称的曲线方程为（\quad）

\choices
    {\( y = \ln (1 + \sqrt{x}) \)}
    {\( y = \ln (1 - \sqrt{x}) \)}
    {\( y = -\ln (1 + \sqrt{x}) \)}
    {\( y = -\ln (1 - \sqrt{x}) \)}
\end{problem}

\begin{answer}
A
\end{answer}

\begin{solution}
原曲线上 \(x\geqslant0\)，且 \(y=\e^{2x}-2\e^x+1=(\e^x-1)^2\). 关于直线 \(y=x\) 对称后交换横、纵坐标，得到 \(x=(\e^y-1)^2\)，并且原曲线的横坐标条件变为反曲线的 \(y\geqslant0\). 因此 \(\e^y-1=\sqrt{x}\)，从而 \(y=\ln(1+\sqrt{x})\)，故选 A.
\end{solution}

\begin{problem}
曲线 \(\frac{x^2}{10 - m} + \frac{y^2}{6 - m} = 1\)（\(m < 6\)）与曲线 \(\frac{x^2}{5 - n} + \frac{y^2}{9 - n} = 1\)（\(5 < n < 9\)）的（\quad）

\choices
    {离心率相等}
    {焦距相等}
    {焦点相同}
    {准线相同}
\end{problem}

\begin{answer}
B
\end{answer}

\begin{solution}
第一条曲线中 \(m<6\)，所以它是椭圆，且半长轴平方与半短轴平方之差为 \((10-m)-(6-m)=4\)，故其焦半距 \(c=2\)，焦距为 \(2c=4\).

第二条曲线中 \(5<n<9\)，改写为 \(\dfrac{y^2}{9-n}-\dfrac{x^2}{n-5}=1\)，它是双曲线，且 \(c^2=(9-n)+(n-5)=4\)，所以同样有 \(c=2\)，焦距为 \(4\). 因此两曲线的焦距相等，故选 B.
\end{solution}

\section{填空题}
\begin{problem}
方程 \(\log_2(x - 1) = 2 - \log_2(x + 1)\) 的解为\fillinblank{}.
\end{problem}

\begin{answer}
\(\sqrt{5}\)
\end{answer}

\begin{solution}
定义域要求 \(x-1>0\)，即 \(x>1\). 在此条件下，原方程化为 \(\log_2(x-1)+\log_2(x+1)=2\)，即 \(\log_2(x^2-1)=2\). 所以 \(x^2-1=4\)，得 \(x=\pm\sqrt5\). 结合 \(x>1\)，舍去负根，故解为 \(x=\sqrt5\).
\end{solution}

\begin{problem}
设 \( g(x)=\begin{cases}
\e^{x},&x\leqslant0,\\
\ln x,&x>0,
\end{cases} \) 则 \( g\left(g\left(\frac{1}{2}\right)\right)= \)\fillinblank{}.
\end{problem}

\begin{answer}
\(\frac{1}{2}\)
\end{answer}

\begin{solution}
由于 \(\dfrac12>0\)，有 \(g\left(\dfrac12\right)=\ln\dfrac12=-\ln2<0\). 因此第二次使用分段函数时取第一段，得到 \(g\left(g\left(\dfrac12\right)\right)=g(-\ln2)=\e^{-\ln2}=\dfrac12\).
\end{solution}

\begin{problem}
如图，半径为 \(2\) 的半球内有一内接正六棱锥 \(P - ABCDEF\)，则此正六棱锥的侧面积是\fillinblank{}.

\bitmapfigure[max width=0.82\linewidth]{img/2006/liaoning_liberal/q15_fig1.png}
\FloatBarrier
\end{problem}

\begin{answer}
\(6\sqrt{7}\)
\end{answer}

\begin{solution}
设半球球心为 \(O\)，正六边形底面的中心为 \(O\)，正六边形的顶点在半径为 \(2\) 的圆周上，正六棱锥顶点为球面上的点 \(P\)，则 \(PO=2\). 设 \(M\) 是正六边形一边的中点，则正六边形边长为 \(2\)，其中心到边的距离为 \(OM=\sqrt3\). 在直角三角形 \(POM\) 中，\(PM=\sqrt{PO^2+OM^2}=\sqrt7\). 六个侧面全等，每个侧面的面积为 \(\dfrac12\cdot2\cdot\sqrt7=\sqrt7\)，所以侧面积为 \(6\sqrt7\).
\end{solution}

\begin{problem}
\(5\) 名乒乓球队员中，有 \(2\) 名老队员和 \(3\) 名新队员. 现从中选出 \(3\) 名队员排成 \(1\)，\(2\)，\(3\) 号参加团体比赛，则入选的 \(3\) 名队员中至少有 \(1\) 名老队员，且 \(1\)，\(2\) 号中至少有 \(1\) 名新队员的排列方法有\fillinblank{} 种.
\end{problem}

\begin{answer}
48
\end{answer}

\begin{solution}
不加限制时，从 \(5\) 名队员中选出并排成三号位，共有 \(\mathrm A_5^3=5\cdot4\cdot3=60\) 种.

不满足“至少有一名老队员”的排列，三人全为新队员，共 \(\mathrm A_3^3=6\) 种；不满足“1，2号中至少有一名新队员”的排列，1，2号均为两名老队员，排法为 \(2!\)，3号从三名新队员中选取，排法为 \(2!\cdot3=6\) 种. 这两类情形不相交，因此所求为 \(60-6-6=48\) 种.
\end{solution}

\section{解答题}
\begin{problem}
已知函数 \(f(x) = \sin^2 x + 2\sin x\cos x + 3\cos^2 x\)，\(x \in \R\). 求：

\begin{enumerate}
\item 函数 \( f(x) \) 的最大值及取得最大值的自变量 \( x \) 的集合；
\item 函数 \( f(x) \) 的单调增区间.
\end{enumerate}
\end{problem}

\begin{solution}
\begin{enumerate}
\item 利用 \(\sin^2x+\cos^2x=1\) 及二倍角公式，有 \(f(x)=2+\cos2x+\sin2x=2+\sqrt2\sin\left(2x+\dfrac{\pi}{4}\right)\). 因此 \(f(x)\leqslant2+\sqrt2\)，当且仅当 \(2x+\dfrac{\pi}{4}=\dfrac{\pi}{2}+2k\pi\)，即 \(x=k\pi+\dfrac{\pi}{8}\)，\(k\in\Z\) 时取得等号. 故最大值为 \(2+\sqrt2\)，取得最大值的自变量集合为 \(\left\{k\pi+\dfrac{\pi}{8}\mid k\in\Z\right\}\).
\item \(f'(x)=2\cos2x-2\sin2x=2\sqrt2\cos\left(2x+\dfrac{\pi}{4}\right)\). 函数在导数非负的区间上单调增加，即 \(-\dfrac{\pi}{2}+2k\pi\leqslant2x+\dfrac{\pi}{4}\leqslant\dfrac{\pi}{2}+2k\pi\). 解得 \(-\dfrac{3\pi}{8}+k\pi\leqslant x\leqslant\dfrac{\pi}{8}+k\pi\)，所以单调增区间为 \(\left[k\pi-\dfrac{3\pi}{8},k\pi+\dfrac{\pi}{8}\right]\)，\(k\in\Z\).
\end{enumerate}
\end{solution}

\begin{problem}
甲，乙两班各派 \(2\) 名同学参加年级数学竞赛，参赛同学成绩及格的概率都为 \(0.6\)，且参赛同学的成绩相互之间没有影响. 求：

\begin{enumerate}
\item 甲，乙两班参赛同学中各有 1 名同学成绩及格的概率；
\item 甲，乙两班参赛同学中至少有 1 名同学成绩及格的概率.
\end{enumerate}
\end{problem}

\begin{solution}
\begin{enumerate}
\item 对甲班，恰有 1 名同学成绩及格的概率为 \(\mathrm C_2^1(0.6)(0.4)=0.48\)，乙班同样为 \(0.48\). 由于两班参赛同学的成绩相互独立，所求概率为 \(0.48^2=0.2304\).
\item 四名同学的成绩都不及格的概率为 \((1-0.6)^4=0.4^4=0.0256\)，所以至少有 1 名同学成绩及格的概率为 \(1-0.0256=0.9744\).
\end{enumerate}
\end{solution}

\begin{problem}
已知正方形 \(ABCD\)，\(E\)，\(F\) 分别是边 \(AB\)，\(CD\) 的中点，将 \(\triangle ADE\) 沿 \(DE\) 折起. 如图所示，记二面角 \(A - DE - C\) 的大小为 \(\theta\)（\(0 < \theta < \pi\)）.

\begin{enumerate}
\item 证明： \(BF \parallel\) 平面 \(ADE\)；
\item 若 \(\triangle ACD\) 为正三角形，试判断点 \(A\) 在平面 \(BCDE\) 内的射影 \(G\) 是否在直线 \(EF\) 上. 说明你的结论，并求角 \(\theta\) 的余弦值.
\end{enumerate}

\FigureLayoutDeclare{img/2006/liaoning_liberal/q19_fig1.png}{14.0cm}{5bp 5bp 5bp 5bp}
\bitmapfigure[max width=0.82\linewidth]{img/2006/liaoning_liberal/q19_fig1.png}
\FloatBarrier
\end{problem}

\begin{solution}
\begin{enumerate}
\item 在正方形中，\(EB\parallel FD\) 且 \(EB=FD\)，所以四边形 \(EBFD\) 是平行四边形，从而 \(BF\parallel ED\). 又 \(ED\subset\) 平面 \(ADE\)，且 \(B\notin\) 平面 \(ADE\)，因此 \(BF\parallel\) 平面 \(ADE\).
\item 设正方形边长为 \(2t\)，其中 \(t>0\). 折叠后 \(AD=CD=2t\)，且 \(E\)，\(F\) 分别为 \(AB\)，\(CD\) 的中点，所以 \(AE=t\)，\(EF=2t\). 由 \(\triangle ACD\) 为正三角形，\(AF\) 是中线，故 \(AF=\sqrt3t\)，于是 \(AE^2+AF^2=EF^2\)，从而 \(AE\perp AF\).

设点 \(A\) 在平面 \(BCDE\) 内的射影为 \(G\). 因为 \(AC=AD\)，所以 \(AG^2+GC^2=AG^2+GD^2\)，即 \(GC=GD\). 因此 \(G\) 在 \(CD\) 的垂直平分线上. 又 \(EF\) 是 \(CD\) 的垂直平分线，故 \(G\in EF\).

在直线 \(EF\) 上建立有向坐标，取 \(E\) 的坐标为 \(0\)，\(F\) 的坐标为 \(2t\)，设 \(G\) 的坐标为 \(u\). 由 \(AG\perp\) 平面 \(BCDE\)，在两个直角三角形中分别有 \(AE^2=AG^2+u^2=t^2\)，\(AF^2=AG^2+(2t-u)^2=3t^2\). 两式相减得 \((2t-u)^2-u^2=2t^2\)，解得 \(u=\dfrac{t}{2}\)，所以 \(G\) 在线段 \(EF\) 上，\(EG=\dfrac{t}{2}\)，\(GF=\dfrac{3t}{2}\)，且 \(AG^2=t^2-u^2=\dfrac{3t^2}{4}\).

在平面 \(BCDE\) 内，\(D\) 到直线 \(EF\) 的距离为 \(t\). 因此 \(\triangle DEG\) 的面积为 \(\dfrac12\cdot\dfrac{t}{2}\cdot t=\dfrac{t^2}{4}\). 令 \(H\) 为 \(G\) 在直线 \(DE\) 上的垂足，则 \(GH\perp DE\). 另一方面，\(DE=t\sqrt5\)，故 \(\dfrac12DE\cdot GH=\dfrac{t^2}{4}\)，从而 \(GH=\dfrac{t}{2\sqrt5}\).

因为 \(AG\perp\) 平面 \(BCDE\)，所以 \(AG\perp GH\)，得 \(AH^2=AG^2+GH^2=\dfrac{3t^2}{4}+\dfrac{t^2}{20}=\dfrac{4t^2}{5}\)，即 \(AH=\dfrac{2t}{\sqrt5}\). 直线 \(AH\) 与 \(HG\) 分别垂直于棱 \(DE\)，所以 \(\angle AHG=\theta\)，于是 \(\cos\theta=\dfrac{HG}{AH}=\dfrac14\).
\end{enumerate}
\end{solution}

\begin{problem}
已知等差数列 \(\{a_{n}\}\) 的前 \(n\) 项和为 \(S_{n} = pn^{2} - 2n + q\)（\(p\in\R\)，\(q\in\R\)），\(n \in \N_{+}\).

\begin{enumerate}
\item 求 \(q\) 的值；
\item 若 \( a_{1} \) 与 \( a_{5} \) 的等差中项为 \(18\)，\( b_{n} \) 满足 \( a_{n} = 2 \log_{2} b_{n} \)，求数列 \(\{b_n\}\) 的前 \(n\) 项和.
\end{enumerate}
\end{problem}

\begin{solution}
\begin{enumerate}
    \item 由等差数列前 \(n\) 项和的性质，\(a_n=S_n-S_{n-1}\). 当 \(n\geqslant2\) 时，
    \[
    a_n=pn^2-2n+q-\left(p(n-1)^2-2(n-1)+q\right)=(2n-1)p-2
    \]
    特别地，\(a_1=S_1=p-2+q\)，\(a_2=3p-2\)，\(a_3=5p-2\). 因为 \(\{a_n\}\) 是等差数列，\(a_2-a_1=a_3-a_2\)，即 \(2p-q=2p\)，所以 \(q=0\).
    \item 由 \(q=0\)，并结合 \(a_1=p-2\) 与上式，得对所有 \(n\in\N_+\)，\(a_n=(2n-1)p-2\). 因 \(\{a_n\}\) 为等差数列，\(a_1\) 与 \(a_5\) 的等差中项就是 \(a_3\)，所以 \(a_3=18\)，即 \(5p-2=18\)，解得 \(p=4\). 于是 \(a_n=8n-6\).

由 \(a_n=2\log_2 b_n\)，得 \(\log_2 b_n=4n-3\)，所以 \(b_n=2^{4n-3}=2\cdot16^{n-1}\). 因此数列 \(\{b_n\}\) 的前 \(n\) 项和为 \(2\dfrac{16^n-1}{16-1}=\dfrac{2}{15}(16^n-1)\).
\end{enumerate}
\end{solution}

\begin{problem}
已知函数 \( f(x) = \frac{1}{3} ax^3 + (a + d)x^2 + (a + 2d)x + d \)，\( g(x) = ax^2 + 2(a + 2d)x + a + 4d \). 其中 \( a > 0 \)，\( d > 0 \). 设 \( x_0 \) 为 \( f(x) \) 的极小值点，\( x_1 \) 为 \( g(x) \) 的极值点，\( g(x_2) = g(x_3) = 0 \)，并且 \( x_2 < x_3 \). 将点 \( (x_0, f(x_0)) \)，\( (x_1, g(x_1)) \)，\( (x_2, 0) \)，\( (x_3, 0) \) 依次记为 \(A\)，\(B\)，\(C\)，\(D\).

\begin{enumerate}
\item 求 \(x_0\) 的值；
\item 若四边形 \(APCD\) 为梯形且面积为 1，求 \(a\)，\(d\) 的值.
\end{enumerate}
\end{problem}

\begin{solution}
\begin{enumerate}
\item 题面第二问将四边形写作 \(APCD\)，但前文定义的四个点依次为 \(A\)，\(B\)，\(C\)，\(D\)；以下按点列及梯形条件，将该四边形理解为 \(ABCD\).

由
\[
f'(x)=ax^2+2(a+d)x+a+2d=a\left(x+1\right)\left(x+1+\frac{2d}{a}\right)
\]
且 \(a>0\)，\(d>0\)，可知 \(f'(x)\) 在 \(x=-1\) 处由负变正，所以 \(f(x)\) 的极小值点为 \(x_0=-1\).
\item 由
\[
g'(x)=2ax+2(a+2d)
\]
得 \(x_1=-1-\dfrac{2d}{a}\). 又
\[
g(x)=a\left(x+1\right)\left(x+1+\frac{4d}{a}\right)
\]
所以 \(x_2=-1-\dfrac{4d}{a}\)，\(x_3=-1\). 直接计算得
\[
f(-1)=-\frac{a}{3}+(a+d)-(a+2d)+d=-\frac{a}{3}
\]
并且
\[
g(x_1)=a\left(-\frac{2d}{a}\right)\left(\frac{2d}{a}\right)=-\frac{4d^2}{a}
\]
因此
\(A\left(-1,-\frac{a}{3}\right)\)，\(B\left(-1-\frac{2d}{a},-\frac{4d^2}{a}\right)\)，\(C\left(-1-\frac{4d}{a},0\right)\)，\(D(-1,0)\).
由于 \(CD\) 水平，而 \(AD\) 为竖直线，且 \(BC\) 不可能竖直，所以梯形 \(ABCD\) 的两底只能是 \(AB\) 与 \(CD\)，即 \(AB\parallel CD\). 于是
\[
-\frac{a}{3}=-\frac{4d^2}{a}
\]
从而 \(a^2=12d^2\). 由 \(a>0\)，\(d>0\)，得 \(a=2\sqrt3\,d\).

此时 \(|AB|=\dfrac{2d}{a}\)，\(|CD|=\dfrac{4d}{a}\)，两底间的距离为 \(\dfrac{a}{3}\)，故梯形面积为
\[
S=\frac12\left(\frac{2d}{a}+\frac{4d}{a}\right)\frac{a}{3}=d
\]
由 \(S=1\) 得 \(d=1\)，从而 \(a=2\sqrt3\).
\end{enumerate}
\end{solution}

\begin{problem}
已知点 \(A(x_{1},y_{1})\)， \(B(x_{2},y_{2})\)（\(x_{1}\neq 0\)，\(x_{2}\neq 0\)）是抛物线 \(y^{2} = 2px\)（\(p > 0\)）上的两个动点. \(O\) 是坐标原点，向量 \(\overrightarrow{OA}\)， \(\overrightarrow{OB}\) 满足 \(\left|\overrightarrow{OA} +\overrightarrow{OB}\right| = \left|\overrightarrow{OA} -\overrightarrow{OB}\right|\)，设圆 \(C\) 的方程为 \(x^{2} + y^{2} - (x_{1} + x_{2})x - (y_{1} + y_{2})y = 0\).

\begin{enumerate}
\item 证明线段 \(AB\) 是圆 \(C\) 的直径；
\item 当圆 \(C\) 的圆心到直线 \(x - 2y = 0\) 的距离的最小值为 \( \frac{2\sqrt{5}}{5} \) 时，求 \(p\) 的值.
\end{enumerate}
\end{problem}

\begin{solution}
\begin{enumerate}
\item 设 \(\bs u=\overrightarrow{OA}\)，\(\bs v=\overrightarrow{OB}\). 由已知
\(\lvert\bs u+\bs v\rvert=\lvert\bs u-\bs v\rvert\)，两边平方并化简得 \(\bs u\cdot\bs v=0\)，即 \(x_1x_2+y_1y_2=0\).

圆 \(C\) 的方程可写为
\[
\left(x-\frac{x_1+x_2}{2}\right)^2+\left(y-\frac{y_1+y_2}{2}\right)^2=\frac{(x_1-x_2)^2+(y_1-y_2)^2}{4}
\]
所以圆心是线段 \(AB\) 的中点. 又由 \(x_1x_2+y_1y_2=0\)，代入点 \(A\) 的坐标可得
\[
x_1^2+y_1^2-(x_1+x_2)x_1-(y_1+y_2)y_1=-(x_1x_2+y_1y_2)=0
\]
故 \(A\) 在圆 \(C\) 上；同理 \(B\) 也在圆 \(C\) 上. 因此 \(AB\) 经过圆心且端点在圆上，线段 \(AB\) 是圆 \(C\) 的直径.
\item 设 \(y_1=u\)，\(y_2=v\). 由抛物线方程，\(x_1=\dfrac{u^2}{2p}\)，\(x_2=\dfrac{v^2}{2p}\). 由于 \(x_1\)，\(x_2\neq0\)，有 \(u\)，\(v\neq0\). 正交条件给出
\[
\frac{u^2v^2}{4p^2}+uv=0
\]
从而 \(uv=-4p^2\). 令 \(s=u+v\)，则
\(x_1+x_2=\frac{u^2+v^2}{2p}=\frac{s^2+8p^2}{2p}=\frac{s^2}{2p}+4p\)，\(y_1+y_2=s\).
圆心到直线 \(x-2y=0\) 的距离为
\[
d(s)=\frac{\left|x_1+x_2-2(y_1+y_2)\right|}{2\sqrt5}=\frac{\left|(s-2p)^2+4p^2\right|}{4p\sqrt5}
\]
其中分子恒为正，且 \(s\) 可取任意实数，所以当 \(s=2p\) 时取得最小值，且
\[
d_{\min}=\frac{4p^2}{4p\sqrt5}=\frac{p}{\sqrt5}
\]
由题设 \(d_{\min}=\dfrac{2\sqrt5}{5}=\dfrac{2}{\sqrt5}\)，得 \(p=2\).
\end{enumerate}
\end{solution}
